% Lösungen Einheit 2 – Prozentsatz berechnen (zusammenbau v0.1)
\einheitenkopf{Einheit 2 · Prozentsatz berechnen}

\erg{5}{a) $20\,\%$ \quad b) $75\,\%$ \quad c) $70\,\%$ \quad d) $50\,\%$}
\erg{6}{a) $17\,\%$ \quad b) $68\,\%$ \quad c) $3\,\%$ \quad d) $45\,\%$ \quad e) $81\,\%$ \quad f) $\frac{19}{50} = \frac{38}{100} = 38\,\%$ \quad g) $\frac{7}{25} = \frac{28}{100} = 28\,\%$ \quad h) $\frac{13}{20} = \frac{65}{100} = 65\,\%$ \quad i) $27 : 60 = 0{,}45 = 45\,\%$ \quad j) $153 : 360 = 0{,}425 = 42{,}5\,\%$ \quad k) $91 : 140 = 0{,}65 = 65\,\%$}
\erg{7}{a) $17 : 24 = 0{,}7083\ldots \approx 70{,}8\,\%$ \quad b) $5 : 7 = 0{,}7142\ldots \approx 71{,}4\,\%$ \quad c) $38 : 55 = 0{,}6909\ldots \approx 69{,}1\,\%$ \quad d) Ganzes $20$; $6 : 20 = 0{,}3 = 30\,\%$ \quad e) Ganzes $25$ Spiele; $18 : 25 = 0{,}72 = 72\,\%$ \quad f) Ganzes $120$ Brote; $72 : 120 = 0{,}6 = 60\,\%$ \quad g) Rabatt $20\,€$; $20 : 80 = 0{,}25 = 25\,\%$ \quad h) Rabatt $9\,€$; $9 : 45 = 0{,}2 = 20\,\%$ \quad i) Rabatt $54\,€$; $54 : 360 = 0{,}15 = 15\,\%$}
\erg{8}{a) Ganzes $24$ Lose; $3 : 24 = 0{,}125 = 12{,}5\,\%$ \quad b) Ganzes $15$; $4 : 15 \approx 0{,}267 \approx 26{,}7\,\%$ \quad c) $57\,000 : 480\,000 = 0{,}11875 \approx 11{,}9\,\%$ \quad d) $136 : 380 \approx 0{,}358 \approx 35{,}8\,\%$ \quad e) abgelehnt $64 - 47 = 17$; $17 : 64 \approx 0{,}266 \approx 26{,}6\,\%$ \quad f) verspätet $92 - 79 = 13$; $13 : 92 \approx 0{,}141 \approx 14{,}1\,\%$ \quad g) $1\,080 : 40\,000 = 0{,}027 = 2{,}7\,\%$ \quad h) $175 : 12\,500 = 0{,}014 = 1{,}4\,\%$ \quad i) $6$ von $13$ sind größer; $6 : 13 \approx 0{,}462 \approx 46{,}2\,\%$ \quad j) $5$ von $9$ warteten länger; $5 : 9 \approx 0{,}556 \approx 55{,}6\,\%$}
\erg{9}{a) Fehler: Ganzes durch Teil gerechnet. Richtig: $12 : 48 = 0{,}25 = 25\,\%$. \quad b) Fehler: Ganzes durch Teil gerechnet (und die Dezimalzahl als Prozentzahl gelesen). Richtig: $18 : 45 = 0{,}4 = 40\,\%$. \quad c) Fehler: Lea teilt durch den Rest (die Äpfel), nicht durch das Ganze. Richtig: Ganzes $24$; $6 : 24 = 0{,}25 = 25\,\%$. \quad d) Gefragt ist, welcher Anteil vom Ganzen der Teil ist – das Ganze sind die $100\,\%$, also steht es im Nenner. \quad e) Ja: Beide Male ist der Teil die Hälfte des Ganzen, also $50\,\%$. \quad f) Nein: Ein Teil ist höchstens so groß wie das Ganze, also höchstens $100\,\%$.}
\erg{10}{a) $25\,\%$ ($1$ von $4$ Teilen) \quad b) $70\,\%$ ($7$ von $10$ Teilen) \quad c) $1$ Teilnehmer: $\frac{100}{120}\,\%$; $54$ Teilnehmer: $54 \cdot \frac{100}{120}\,\% = 45\,\%$ \quad d) übrig $3$ GB, $3 : 20 = 0{,}15 = 15\,\%$; pro Tag $17 : 20 = 0{,}85$ GB, für $10$ Tage bräuchte man $10 \cdot 0{,}85 = 8{,}5$ GB – es reicht nicht. \quad e) Ganzes $27$; $12 : 27 \approx 0{,}444 \approx 44{,}4\,\%$ – weniger als die Hälfte: Emma hat die meisten Stimmen, aber nicht die Mehrheit. \quad f) Montag $36 : 240 = 0{,}15 = 15\,\%$; Dienstag $27 : 150 = 0{,}18 = 18\,\%$ – am Dienstag.}
